\chapter{Boxes}

\section{Boxing examples}

The numbering of the various boxes (such as Definitions, Theorems, Exercises, etc.)\ is automatically reset at the beginning of each chapter. This means that every time a new chapter starts, the box counter resets to 1.

Additionally, the numbering between different types of boxes (such as Definitions, Theorems, Exercises, etc.)\ is independent. This means that the numbering for Definitions will not affect the numbering for Theorems, and so on.

\begin{definition}{Definition term}
    Example of a colored ``Definition'' box.
\end{definition}

\begin{theorem}{Theorem name}
    Example of a colored ``Theorem'' box.
\end{theorem}

\begin{corollary}{Corollary name}
    Example of a colored ``Corollary'' box.
\end{corollary}

\begin{exercise}{Exercise title}
    Example of a colored ``Exercise'' box.
\end{exercise}

\begin{observation}{Observation title}
    Example of a colored ``Observation'' box.
\end{observation}

The source code is the same for both Italian and English documents:
\begin{quote}
\begin{verbatim}
\begin{theorem}{Theorem name}
    Example of a colored ``Theorem'' box.
\end{theorem}
\end{verbatim}
\end{quote}
